
Performance of the \textbf{heterogeneous ensemble} on the phone-captured dataset (see Table~\ref{tab:hetero_phone}) exhibits a substantial drop relative to clean validation results, reflecting the severity of real-world distribution shift.
Introducing architectural diversity improves realized ensemble accuracy compared to homogeneous CNN aggregation.
More importantly, oracle accuracy increases to 72.49\%, substantially exceeding realized ensemble performance.
The magnitude of this oracle gap is notably larger than on clean data.
This indicates that under real-world shift, different experts succeed on different subsets of the data.
In other words, complementary expertise becomes more pronounced as the distribution diverges from training conditions.
Despite this increased diversity, uniform averaging fails to recover oracle performance.
The ensemble remains significantly below its theoretical upper bound.
This suggests that the heterogeneous ensemble is not capacity-limited, but limited by its aggregation strategy at inference time.

The presence of substantial unrealized oracle potential motivates the introduction of an adaptive routing mechanism, which aims to exploit input-dependent specialization among experts.
